% Revision/source crosswalk: analyses/018-2026-09-08-results-materials/observations/O014-manuscript-polish.md
% Evidence: Analysis 018 observations O002/O003/O007/O011; no new measurements.
\begin{abstract}
Activation sparsity can reduce inference cost, but the value of a
sparsification intervention depends on where it acts and what accompanies it.
We study pressure toward zero, thresholding and site placement through
54 matched pretraining conditions on MiniPile, using Pythia-14M, 70M and
410M. We measure validation loss, activation
distributions and model-wide sparsity: the fraction of counted scalar multiplications
with a zero activation operand. Pressure's added value changes with threshold
and target sites. A seven-site recipe can yield greater model-wide sparsity
than a four-site recipe despite fewer feed-forward (FFN) zeros, because it also sparsifies
internal attention products. At the largest tested threshold, the seven-site
recipe with pressure improves both loss and sparsity over its four-site
counterpart at all three sizes. On one RTX5090, a selected fused implementation
achieves a $1.234\times$ geometric mean full-model speedup across the 30
14M checkpoints for uncached, 2,048-token inference. Sparse paths add only
$1.043\times$ over the same fusion on average, and attention skipping slows
every checkpoint. By connecting matched intervention effects to activation
distributions, affected computation and measured execution, the study provides
a basis for choosing where and how to induce useful sparsity.
\end{abstract}
